\documentclass{article}
\usepackage[T1]{fontenc}
\usepackage{lmodern}
\usepackage{graphicx}
\usepackage{amsmath,amssymb}
\usepackage[a4paper,margin=2cm]{geometry}
\usepackage[absolute,overlay]{textpos}
\setlength{\TPHorizModule}{1mm}
\setlength{\TPVertModule}{1mm}
\usepackage[indonesian]{babel}
\usepackage{hyperref}
\setlength{\emergencystretch}{2em}
% unit: Halaman 23

\begin{document}

% === Halaman 23 (python/native) ===

23

• Selanjutnya, 


Rtate mungkin adalah state, 


atthattMZE mungkin adalah atthattime, 


heVe mungkin adalah here. 

• Jadi, kita memperoleh pemetaan baru:


 R → s


 M → i


 Z → m


 V → r

\end{document}
